\subsection{CONSTRUCTION OF THE FREE LIE ALGEBRA}

DEFINITION 1. The free Lie algebra over the set \(\mathbf{X}\) is the quotient algebra

\[
\mathrm{L} (\mathrm{X}) = \operatorname{Lib} (\mathrm{X}) / \mathrm{a},
\]

where a is the two-sided ideal of \(\operatorname{Lib}(\mathbf{X})\) generated by the elements of one of the forms

\begin{enumerate}
\def\labelenumi{(\arabic{enumi})}
\tightlist
\item
\end{enumerate}

\[
\mathrm{Q} (a) = a. a \text {   for   } a \text {   in   } \operatorname{Lib} (\mathrm{X}),\tag{2}
\]

\[
J (a, b, c) = a. (b. c) + b. (c. a) + c. (a. b)
\]

for \(a, b, c\) in \(\operatorname{Lib}(\mathbf{X})\) .

Clearly \(\mathrm{L}(\mathrm{X})\) is a Lie K-algebra; the composition of two elements \(u, v\) of \(\mathrm{L}(\mathrm{X})\) will be denoted by \([u, v]\) . When it is necessary to indicate the ring \(\mathbf{K}\) , we write \(\mathrm{L}_{\mathbb{K}}(\mathbf{X})\) for \(\mathrm{L}(\mathbf{X})\) .

The following proposition justifies the name free Lie algebra given to L(X).

PROPOSITION 1. Let \(\psi\) be the canonical mapping of \(\operatorname{Lib}(\mathbf{X})\) onto \(\operatorname{L}(\mathbf{X})\) and \(\phi\) the restriction of \(\psi\) to \(\mathbf{X}\) . For every mapping \(f\) of \(\mathbf{X}\) into a Lie algebra \(\mathfrak{g}\) , there exists one and only one homomorphism \(\mathrm{F}:\mathrm{L}(\mathrm{X})\to \mathfrak{g}\) such that \(f = \mathrm{F}\circ \phi\) .

\begin{enumerate}
\def\labelenumi{(\alph{enumi})}
\item
  Existence of F: let h be the homomorphism of \(\operatorname{Lib}(\mathbf{X})\) into g extending \(f\) (no. 1) . For all a in \(\operatorname{Lib}(\mathbf{X})\) , \(h(\mathbf{Q}(a)) = h(a.a) = [h(a), h(a)] = 0\) ; similarly, the Jacobi identity satisfied by g implies that \(h(\mathbf{J}(a, b, c)) = 0\) for a, b, c in \(\operatorname{Lib}(\mathbf{X})\) . It follows that \(h(a) = 0\) , whence there is a homomorphism F of \(\operatorname{L}(\mathbf{X})\) into g such that \(h = F \circ \phi\) . By restricting to X, we obtain \(f = F \circ \phi\) .
\item
  Uniqueness of F: let \(\mathrm{F}^{\prime}:\mathrm{L}(\mathrm{X})\to\mathfrak{g}\) be a homomorphism such that \(f=\mathrm{F}^{\prime}\circ\phi\) . The homomorphisms \(F\circ\psi\) and \(F^{\prime}\circ\psi\) of \(\operatorname{Lib}(X)\) into g coincide on X and hence are equal; as \(\psi\) is surjective, \(F=F^{\prime}\) .
\end{enumerate}

COROLLARY 1. The family \((\phi (x))_{x\in \mathbf{X}}\) is free over \(\mathbf{K}\) in \(\mathrm{L}(\mathbf{X})\) .

Let \(x_{1}, x_{2}, \ldots, x_{n}\) be distinct elements of X and \(\lambda_{1}, \ldots, \lambda_{n}\) be elements of K such that

\[
\lambda_ {1}. \phi (x _ {1}) + \dots + \lambda_ {n}. \phi (x _ {n}) = 0.\tag{3}
\]

Let \(\mathfrak{g}\) be the commutative Lie algebra with \(\mathbf{K}\) as underlying module. For \(i = 1,2,\ldots ,n\) , there exists a homomorphism \(\mathrm{F}_i\) of \(\mathbf{L}(\mathbf{X})\) into \(\mathfrak{g}\) such that \(\mathrm{F_i}(\phi (x_i)) = 1\) and \(\mathrm{F_i}(\phi (x)) = 0\) for \(x\neq x_{i}\) (Proposition 1); applying \(\mathrm{F}_i\) to relation (3), we obtain \(\lambda_{i} = 0\) .

COROLLARY 2. Let \(\mathfrak{a}\) be a Lie algebra. Every extension of \(\mathrm{L}(\mathrm{X})\) by \(\mathfrak{a}\) is inessential.

Let \(\mathfrak{a} \xrightarrow{\lambda} \mathfrak{g} \xrightarrow{\mu} \mathrm{L}(\mathrm{X})\) be such an extension (Chapter I, § 1, no. 7). As \(\mu\) is surjective, there exists a mapping \(f\) of \(\mathrm{X}\) into \(\mathfrak{g}\) such that \(\phi = \mu \circ f\) . Let \(\mathrm{F}\) be the homomorphism of \(\mathrm{L}(\mathrm{X})\) into \(\mathfrak{g}\) such that \(f = \mathrm{F} \circ \phi\) (Proposition 1). Then \((\mu \circ \mathrm{F}) \circ \phi = \mu \circ f = \phi\) and Proposition 1 shows that \(\mu \circ \mathrm{F}\) is the identity automorphism of \(\mathrm{L}(\mathrm{X})\) . The given extension is therefore inessential (Chapter I, § 1, no. 7, Proposition 6 and Definition 6).

As the ring K is non-zero, Corollary 1 to Proposition 1 shows that \(\phi\) is injective. Hence the set X can be identified by means of \(\phi\) with its image in L(X); with this convention, X generates L(X) and every mapping of X into a Lie algebra \(\mathfrak{g}\) can be extended to a Lie algebra homomorphism of L(X) into \(\mathfrak{g}\) .

Remark. When X is empty, M(X) is empty and hence L(X) = \{0\}. If X consists of a single element x, the submodule K.x of L(X) is a Lie subalgebra of

\(L(X)\) ; as X generates \(L(X)\) , Corollary 1 to Proposition 1 shows that \(L(X)\) is a free module with basis \(\{x\}\) .

